% !TEX root = ../../main.tex
% C1.1 — Bối cảnh và lý do chọn đề tài
% Nội dung cần có: Camera giám sát phổ biến, dữ liệu lớn; tìm người thủ công tốn thời gian; người dùng chỉ có ảnh, mô tả hoặc đặc điểm trang phục.
\section{Lý do chọn đề tài}
\label{sec:1.1}

Camera giám sát đang được triển khai ngày càng rộng rãi tại Việt Nam. Theo số liệu của Tổng cục Hải quan được Bộ Thông tin và Truyền thông dẫn lại, trong năm năm gần đây đã có hơn 16 triệu thiết bị camera giám sát được nhập khẩu vào Việt Nam, trung bình khoảng 3,2 triệu thiết bị mỗi năm, và ước tính đến năm 2025 cả nước có hơn 20 triệu camera giám sát đang được sử dụng~\cite{vnexpress2024camera}. Số lượng camera lớn cùng thời gian ghi hình liên tục tạo ra một khối lượng dữ liệu video rất lớn.

Giá trị của hệ thống camera không chỉ nằm ở khả năng răn đe. Tổng hợp bốn mươi năm nghiên cứu cho thấy camera giám sát giúp giảm tội phạm ở mức vừa phải~\cite{piza2019cctv}, nhưng vai trò nổi bật hơn của chúng là cung cấp bằng chứng cho việc truy vết sau sự việc. Một nghiên cứu trên 251.195 vụ việc do Cảnh sát Giao thông Đường sắt Anh ghi nhận cho thấy hình ảnh camera hữu ích cho quá trình điều tra trong khoảng 29\% số vụ, và khi có hình ảnh hữu ích, tỉ lệ vụ việc được giải quyết tăng từ khoảng 20\% lên khoảng 50\%~\cite{ashby2017cctv}. Như vậy, giá trị khai thác thực tế của dữ liệu camera phụ thuộc lớn vào việc có tìm được đúng đoạn hình ảnh cần thiết hay không.

Trên thực tế, nhiều tổ chức vận hành camera ghi hình liên tục và có nhu cầu tìm lại hình ảnh, thời điểm và vị trí xuất hiện của một người cụ thể. Nếu xem lại thủ công từng đoạn video của từng camera, việc tìm kiếm sẽ tốn nhiều thời gian và dễ bỏ sót. Với hệ thống camera trải rộng trên nhiều khu vực, việc tìm kiếm thường được phân công theo khu vực, còn kết quả cần được ghi nhận thành hồ sơ để cấp quản lý theo dõi, đặt ra thêm yêu cầu về phân quyền truy cập dữ liệu và quản lý kết quả.

Bên cạnh đó, thông tin ban đầu về người cần tìm thường không đầy đủ và tồn tại dưới nhiều dạng khác nhau: một hình ảnh tham chiếu, một đoạn mô tả bằng lời, hoặc một vài đặc điểm ngoại hình như giới tính, loại và màu trang phục hay vật mang theo. Vì vậy, một hệ thống tìm kiếm hữu ích cần tiếp nhận được các dạng mô tả đa phương thức này và đưa chúng về cùng một cơ chế so khớp với dữ liệu đã được phân tích từ camera, thay vì chỉ hỗ trợ một hình thức truy vấn duy nhất.

Xuất phát từ những nhu cầu thực tiễn trên, đề tài ``Phát triển hệ thống tìm kiếm người dựa trên mô tả đa phương thức'' được lựa chọn thực hiện.
